% =========================================================================
% PROJECT WALKTHROUGH (companion document to THESIS.tex)
%
% A self-contained explainer of what the project does, the math behind it,
% and how to defend each component. Compile on Overleaf with pdfLaTeX.
% =========================================================================
\documentclass[11pt,a4paper]{article}

% --- Layout ---
\usepackage[margin=2.2cm]{geometry}
\usepackage{setspace}
\linespread{1.15}

% --- Encoding & fonts ---
\usepackage[T1]{fontenc}
\usepackage[utf8]{inputenc}
\usepackage{lmodern}
\usepackage{microtype}

% --- Math ---
\usepackage{amsmath, amssymb, amsthm}
\usepackage{bm}
\usepackage{mathtools}

% --- Tables, lists, graphics, captions ---
\usepackage{booktabs, array, tabularx}
\usepackage{enumitem}
\setlist{topsep=2pt, itemsep=2pt, parsep=0pt}
\usepackage{graphicx}
\usepackage{xcolor}
\usepackage[font=small,labelfont=bf]{caption}

% --- Code listings ---
\usepackage{listings}
\lstdefinestyle{pystyle}{
  language=Python,
  basicstyle=\ttfamily\footnotesize,
  keywordstyle=\color{blue!70!black}\bfseries,
  commentstyle=\color{green!50!black}\itshape,
  stringstyle=\color{red!70!black},
  showstringspaces=false,
  breaklines=true,
  frame=single,
  framesep=4pt,
  backgroundcolor=\color{gray!5},
  numbers=left, numberstyle=\tiny\color{gray}, numbersep=5pt,
}
\lstset{style=pystyle}

% --- Section spacing ---
\usepackage{titlesec}
\titlespacing*{\section}{0pt}{14pt}{6pt}
\titlespacing*{\subsection}{0pt}{10pt}{4pt}

% --- Highlight boxes ---
\usepackage{tcolorbox}
\tcbuselibrary{breakable, skins}
\newtcolorbox{keybox}[1][]{enhanced, breakable,
  colback=blue!4, colframe=blue!50!black, boxrule=0.4pt,
  title=#1, fonttitle=\bfseries,
  left=6pt, right=6pt, top=4pt, bottom=4pt}
\newtcolorbox{notebox}[1][]{enhanced, breakable,
  colback=yellow!8, colframe=orange!60!black, boxrule=0.4pt,
  title=#1, fonttitle=\bfseries,
  left=6pt, right=6pt, top=4pt, bottom=4pt}

% --- Hyperref ---
\usepackage{hyperref}
\hypersetup{colorlinks=true, linkcolor=blue!50!black,
            urlcolor=blue!60!black, citecolor=teal!60!black}

% =========================================================================
\title{\textbf{Project Walkthrough}\\[0.3em]
       \large Self-Supervised Reconstruction-Guided\\
       Edge Pruning for GNN-based IoT Intrusion Detection\\[0.3em]
       \normalsize A companion explainer to \texttt{THESIS.tex}}
\author{Pulkit Mittal}
\date{\today}

\begin{document}
\maketitle

\begin{keybox}[How to read this document]
This is a learning-and-defence companion to the formal thesis
(\texttt{THESIS.tex}). The thesis is written in the voice of a finished
report; this document is written in the voice of \emph{explaining the
project to you}, with extra hand-holding on the math and intuition.
\begin{itemize}
  \item Sections 1--2 explain \emph{what we did} in plain terms.
  \item Sections 3--6 build up the math step by step.
  \item Section 7 explains \emph{why each piece is necessary} via
        ablations.
  \item Section 8 connects math to the empirical numbers.
  \item Section 9 is the elevator-pitch you can memorise for the
        defence.
  \item Section 10 is a one-page cheat-sheet for quick reference.
\end{itemize}
\end{keybox}

\tableofcontents

% =========================================================================
\section{What this project does (plain English)}

We take the standard GNN-based IoT intrusion-detection architecture
from Villegas-Ch et al.\ (IEEE Access), which is a 3-layer Graph
Convolutional Network on the IoT communication graph. We \emph{add two
components on top of it}:
\begin{enumerate}
  \item A small \textbf{autoencoder} trained on normal-class nodes
        only.
  \item A \textbf{differentiable edge pruner} that uses the
        autoencoder's reconstruction error to decide which edges to
        keep.
\end{enumerate}
The output is a \emph{sparsified} graph that the same 3-layer GCN
classifies. We removed about half the edges at zero F1 loss, and the
inference is 1.37$\times$ faster on graphs with $\sim$50\,000 nodes.

\paragraph{What \emph{we} contribute methodologically.}
The original GNN paper trains a classifier. We don't change it. We
change \emph{how the edges feeding into that classifier are selected}.
Specifically, the autoencoder is trained on \textbf{only the normal
class}, so its reconstruction error becomes an
\emph{out-of-distribution score} that the pruner can consume
\emph{without ever seeing attack labels}. This is what makes our pruning
\textbf{label-efficient} --- it works the same whether you have 100\%
of the attack labels or only 5\%.

\paragraph{What "self-supervised" means here.}
The autoencoder's training task --- reconstruct its own input ---
requires no labels. We additionally filter \emph{which} nodes are
allowed to enter the reconstruction loss (normal nodes only). The
attack-label information is therefore never injected into the
autoencoder's gradient. This is precisely the property prior pruning
methods (NeuralSparse, PTDNet) lack.

% =========================================================================
\section{The big picture in one diagram}

\begin{center}
\renewcommand{\arraystretch}{1.25}
\begin{tabular}{p{0.7\textwidth}}
\toprule
\textbf{1.} Raw IoT graph $\mathcal{G} = (V, E, X, A)$ \\
\quad $\downarrow$ \\
\textbf{2.} Autoencoder $\to$ per-node MSE score (trained on normals)\\
\quad $\downarrow$ \\
\textbf{3.} Pruner $\to$ binary keep/prune decision per edge\\
\quad $\downarrow$ \\
\textbf{4.} Sparsified graph $\mathcal{G}'$ with $\sim$50\% fewer edges\\
\quad $\downarrow$ \\
\textbf{5.} 3-layer GCN classifier (same as reference paper)\\
\quad $\downarrow$ \\
\textbf{6.} Per-node predictions $\hat{y}$ (normal / attack)\\
\bottomrule
\end{tabular}
\end{center}

Steps 1, 5, 6 are the standard GNN-IDS pipeline. Steps 2--4 are our
addition. All five trainable components --- autoencoder, pruner, three
GCN layers --- are trained jointly through one loss, end-to-end.

% =========================================================================
\section{Notation and the standard baseline}
\label{sec:notation}

\subsection{The data}

\begin{tabular}{ll}
\toprule
Symbol & Meaning \\
\midrule
$\mathcal{G} = (V, E, X, A)$ & The IoT communication graph \\
$|V| = N$ & Number of nodes (devices) \\
$|E| = M$ & Number of edges (communications) \\
$X \in \mathbb{R}^{N \times d_x}$ & Per-node features ($d_x = 21$) \\
$A \in \mathbb{R}^{N\times N}$ & Weighted adjacency \\
$y_i \in \{0,1\}$ & Node label (0=normal, 1=attack) \\
$L_{\text{train}}, L_{\text{val}}, L_{\text{test}}$ & Labelled train/val/test subsets \\
$x_u, x_v$ & Endpoint feature vectors of edge $(u,v)$ \\
$w_{uv}$ & Edge weight \\
$\hat{y} \in [0,1]^{N\times 2}$ & Model output probabilities \\
\bottomrule
\end{tabular}

\subsection{The standard 3-layer GCN (Villegas-Ch et al.)}

Each GCN layer applies symmetric Laplacian smoothing:
\begin{equation}
H^{(l+1)} \;=\; \sigma\!\left(\widetilde{D}^{-1/2}\,\widetilde{A}\,\widetilde{D}^{-1/2}\,
H^{(l)}\,W^{(l)}\right)
\label{eq:gcn}
\end{equation}
where $\widetilde{A} = A + I$ adds self-loops, $\widetilde{D}_{ii} =
\sum_j \widetilde{A}_{ij}$ is the degree matrix, $H^{(0)} = X$, and
$W^{(l)}$ are learnable weights. ReLU acts as $\sigma$.

The pipeline is:
\[
X \xrightarrow{\text{GCN}_1} \mathbb{R}^{N\times 64}
  \xrightarrow{\text{GCN}_2} \mathbb{R}^{N\times 128}
  \xrightarrow{\text{GCN}_3} \mathbb{R}^{N\times 64}
  \xrightarrow{\text{MLP}}  \hat{Y} \in [0,1]^{N\times 2}.
\]

The classification loss is standard cross-entropy on labelled nodes:
\begin{equation}
\mathcal{L}_{\text{CE}}
= \frac{1}{|L_{\text{train}}|}
  \sum_{i\in L_{\text{train}}}
  -\!\big[\,y_i\,\log \hat{y}_{i,1} + (1-y_i)\,\log \hat{y}_{i,0}\,\big].
\label{eq:lce}
\end{equation}

\begin{notebox}[Intuition for the GCN]
Each GCN layer makes every node \emph{average its neighbours' features},
weighted by edge weight and normalised by degrees. When neighbours
share labels (homophily holds), this denoises the prediction. When
neighbours have different labels (cross-class edge), this
\emph{contaminates} the node's representation. Removing cross-class
edges is therefore both an accuracy lever and a speed lever.
\end{notebox}

% =========================================================================
\section{Our addition 1: class-conditional autoencoder}
\label{sec:ae}

\subsection{Architecture}

A two-layer feed-forward encoder and decoder:
\begin{align}
z_i  &= \mathrm{ReLU}\!\left(W_{e2}\,\mathrm{ReLU}(W_{e1}\,x_i)\right)
        \quad &&\text{(encoder)} \\
\hat{x}_i &= W_{d2}\,\mathrm{ReLU}(W_{d1}\,z_i)
        \quad &&\text{(decoder)}
\end{align}
with hidden dim 64 and latent dim 32.

Per-node mean-squared reconstruction error:
\begin{equation}
\mathrm{MSE}_i \;=\; \frac{1}{d_x}\sum_{k=1}^{d_x}(x_{ik} - \hat{x}_{ik})^2.
\label{eq:mse}
\end{equation}

\subsection{The novelty: class-conditional masking}

Let $S \;=\; \{\,i \in L_{\text{train}} \;:\; y_i = 0\,\}$ be the set of
\emph{normal-class training nodes}. The autoencoder's training loss is
restricted to this set:
\begin{equation}
\boxed{\;\mathcal{L}_{\text{recon}}
\;=\; \frac{1}{|S|}\sum_{i \in S}\| x_i - \hat{x}_i \|^2\;}
\label{eq:lrecon}
\end{equation}

Equation~\eqref{eq:lrecon} is the \emph{only} novelty-bearing line in
the entire method. Everything else --- the AE topology, the Gumbel
trick below, the GCN backbone --- is off-the-shelf.

\subsection{Why the MSE becomes an OOD score}

If the autoencoder has sufficient capacity and is trained only on
normal-class samples, it learns the conditional density
$p_X(\,\cdot \mid y=0\,)$. After training:
\begin{align*}
\mathbb{E}[\,\mathrm{MSE}_i \mid y_i = 0\,] &\;\approx\; 0\\
\mathbb{E}[\,\mathrm{MSE}_i \mid y_i = 1\,] &\;\gg\; 0
\end{align*}
because attack-class samples are \emph{out of the AE's training
distribution} and therefore reconstructed badly.

For an edge $(u, v)$, define
\(\Delta_{uv} = |\mathrm{MSE}_u - \mathrm{MSE}_v|\).
The conditional means of $\Delta_{uv}$ then satisfy:
\begin{align*}
\mathbb{E}[\,\Delta_{uv} \mid (u,v) \text{ same-class}] &\;\approx\; \text{small}\\
\mathbb{E}[\,\Delta_{uv} \mid (u,v) \text{ cross-class}] &\;\gg\; \text{small}.
\end{align*}
So $\Delta_{uv}$ is the discriminative feature the pruner needs to
identify cross-class (noisy) edges, and it is constructed without
referencing the attack labels at any point.

\begin{notebox}[What if we trained the AE on \emph{all} nodes?]
This is the ``full-AE ablation'' in the thesis. The AE then learns the
mixture distribution
$p_X(\,\cdot\,) = \pi\,p_X(\,\cdot\mid y=0) + (1-\pi)\,p_X(\,\cdot\mid y=1)$.
Attack nodes are now in-distribution; their MSE drops to a mid-range
value; $\Delta_{uv}$ no longer discriminates cleanly. The empirical
consequence (Section~\ref{sec:results-link}) is that the selectivity
ratio falls from $1.42\times$ (class-conditional) to $0.88\times$
(full AE), which is \emph{below random pruning}. This is the cleanest
empirical evidence that the class-conditional masking is the necessary
ingredient.
\end{notebox}

% =========================================================================
\section{Our addition 2: differentiable Gumbel-Softmax pruner}
\label{sec:pruner}

\subsection{Edge features and the scorer MLP}

For each edge $(u,v) \in E$, build a 4-dim feature vector:
\begin{equation}
\phi_{uv} \;=\; \big[\,\mathrm{MSE}_u,\; \mathrm{MSE}_v,\;
                       \cos(x_u, x_v),\; w_{uv}\,\big]
\label{eq:phi}
\end{equation}
where $\cos(x_u, x_v) = (x_u^\top x_v)/(\|x_u\|\|x_v\|)$ is cosine
similarity. A small MLP $g_\psi : \mathbb{R}^4 \to \mathbb{R}^2$ outputs
\emph{2-class logits} per edge:
\begin{equation}
\ell_{uv} \;=\; g_\psi(\phi_{uv})
\;=\; \big(\,\ell_{uv,0},\; \ell_{uv,1}\,\big),
\label{eq:logits}
\end{equation}
where index 0 is ``prune'' and index 1 is ``keep''.

\subsection{The differentiability problem}

If we just took the arg-max of $\ell_{uv}$, we would get a binary
keep/prune decision --- but the operation $\arg\max$ has zero gradient
almost everywhere, so we cannot back-propagate through it. We need a
differentiable surrogate. The Gumbel-Softmax with the
Straight-Through estimator solves this.

\subsection{Gumbel-Softmax}

\paragraph{Step 1: Gumbel-distributed noise.}
Sample
\(g_{uv,c} \sim \mathrm{Gumbel}(0,1)\) for $c \in \{0,1\}$. A clean
reparameterisation:
\[g_{uv,c} \;=\; -\log(-\log u),\quad u \sim \mathrm{Uniform}(0,1).\]

\paragraph{Step 2: Soft mask via tempered softmax.}
\begin{equation}
y^{\text{soft}}_{uv} \;=\; \mathrm{softmax}\!\Bigg(\frac{\ell_{uv} + g_{uv}}{\tau}\Bigg)
\;\in\; [0,1]^{2}.
\label{eq:ysoft}
\end{equation}
When $\tau \to 0$ this converges to the argmax distribution; when
$\tau$ is large, it is smooth. We schedule $\tau$ from 1.0 (smooth) to
0.1 (sharp) over training.

\paragraph{Step 3: Hard mask via argmax.}
\begin{equation}
y^{\text{hard}}_{uv} \;=\; \mathrm{onehot}\!\big(\arg\max y^{\text{soft}}_{uv}\big).
\label{eq:yhard}
\end{equation}

\paragraph{Step 4: Straight-Through identity.}
We want \emph{hard forward, soft gradient}. The trick:
\begin{equation}
m_{uv} \;=\; y^{\text{hard}}_{uv,1}
\;+\; y^{\text{soft}}_{uv,1}
\;-\; \mathrm{sg}(y^{\text{soft}}_{uv,1}),
\label{eq:ste}
\end{equation}
where $\mathrm{sg}(\cdot)$ is the stop-gradient operator.

\begin{keybox}[Why the ST identity gives "hard forward / soft gradient"]
\textbf{Forward}: numerically,
$y^{\text{soft}}_{uv,1} - \mathrm{sg}(y^{\text{soft}}_{uv,1}) = 0$,
so $m_{uv} = y^{\text{hard}}_{uv,1}$, which is 0 or 1. \emph{The GCN
truly receives a binary mask.}\\
\textbf{Backward}: $\partial m_{uv}/\partial \ell_{uv}
= \partial y^{\text{soft}}_{uv,1}/\partial \ell_{uv}$, because
$y^{\text{hard}}$ and $\mathrm{sg}(\cdot)$ are constants under
differentiation. \emph{The pruner receives a smooth gradient.}
\end{keybox}

Self-loops are forced kept in both branches by setting $m_{uu} = 1$
after Eq.~\eqref{eq:ste}, to prevent disconnected nodes from producing
NaN.

\subsection{The sparsity regulariser}

To control \emph{how many} edges we keep, we add a quadratic penalty
on the average soft keep-probability:
\begin{equation}
\mathcal{L}_{\text{sparse}}
\;=\; \Bigg(\,\frac{1}{|E\setminus\text{self-loops}|}
                \sum_{(u,v)} y^{\text{soft}}_{uv,1} \;-\; (1-\rho)\Bigg)^{\!2}
\label{eq:lsparse}
\end{equation}
where $\rho \in (0,1)$ is the target sparsity. We use $\rho = 0.5$.
$\mathcal{L}_{\text{sparse}}$ uses no labels of any kind.

% =========================================================================
\section{The tri-partite loss and the supervision-decoupling property}

\subsection{The total loss}

\begin{equation}
\boxed{\;\mathcal{L}_{\text{total}}
\;=\; \underbrace{\mathcal{L}_{\text{CE}}}_{\text{supervised}}
\;+\; \underbrace{\lambda_1\,\mathcal{L}_{\text{recon}}}_{\text{self-supervised}}
\;+\; \underbrace{\lambda_2\,\mathcal{L}_{\text{sparse}}}_{\text{label-free}}\;}
\label{eq:total}
\end{equation}

with $\lambda_1 = 0.1$, $\lambda_2 = 1.0$.

\subsection{Why this is the key property}

\begin{table}[h]
\centering
\caption{Each loss term has a different supervision mode. The
autoencoder's loss path uses no attack labels.}
\begin{tabular}{lllc}
\toprule
\textbf{Term} & \textbf{Affects} & \textbf{Supervision} & \textbf{Uses attack labels?} \\
\midrule
$\mathcal{L}_{\text{CE}}$ & GCN classifier $\theta$ & Supervised on $L_{\text{train}}$ & yes \\
$\mathcal{L}_{\text{recon}}$ & Autoencoder $W_e, W_d$ & Filtered by normal label & \textbf{no} \\
$\mathcal{L}_{\text{sparse}}$ & Pruner $\psi$ & Label-free target & no \\
\bottomrule
\end{tabular}
\end{table}

The pruner $\psi$ receives gradients from \emph{only}
$\mathcal{L}_{\text{CE}}$ (back-propagated through the masked GCN) and
$\mathcal{L}_{\text{sparse}}$. The autoencoder receives gradients
from \emph{only} $\mathcal{L}_{\text{recon}}$ --- which never references
the attack label. Therefore the AE's output (the $\mathrm{MSE}_i$
score) is \emph{invariant to how many attack labels are available}.

That is the property that gives us the flat label-efficiency curve in
the empirical results.

% =========================================================================
\section{Training and inference}

\subsection{Two-phase training schedule}

\paragraph{Phase 1 --- Warm-up (epochs 1--50).}
Set $\lambda_2 = 0$ and $\tau = 1$. No sparsity pressure, soft mask.
This lets the autoencoder converge first --- otherwise the pruner
would start dropping edges before the AE knows which edges are noisy.

\paragraph{Phase 2 --- Pruning (epochs 51--300).}
Enable $\lambda_2 = 1.0$ and geometrically anneal the Gumbel
temperature:
\begin{equation}
\tau_e \;=\; \tau_{\text{start}}^{\,1-t}\;\tau_{\text{end}}^{\,t},
\qquad t = (e-50) / (300-50).
\end{equation}
$\tau$ goes from $1.0$ at epoch 50 to $0.1$ at epoch 300. Annealing
avoids the model getting stuck in the soft regime forever.

\paragraph{Optimiser.} Adam, learning rate $10^{-3}$, weight decay
$5\times 10^{-4}$, gradient clip at norm 5.

\subsection{Inference (deployment mode)}

At test time we drop the Gumbel noise and use deterministic top-k
selection. For each edge compute a score:
\[
s_{uv} \;=\; \ell_{uv,1} - \ell_{uv,0}.
\]
Sort and keep the top $(1-\rho)\,|E_{\setminus\text{self-loop}}|$
edges by $s_{uv}$.

\paragraph{Deployment optimisation.}
For a static IoT topology, run the AE + pruner once to obtain the
fixed pruned graph $\mathcal{G}'$. Thereafter, deployment runs only
the 3-layer GCN backbone on $\mathcal{G}'$. The AE and pruner are
training-time scaffolding only. This is the
\emph{deploy-mode inference} that yields the $1.37\times$ speed-up at
50K-node graphs.

% =========================================================================
\section{Why each piece is necessary (the ablation logic)}

The ablation table in the thesis is a step-by-step \emph{stripping
experiment}: each row removes one design element from the row below
to see how performance degrades.

\begin{table}[h]
\centering
\caption{Each row removes one design element from the row below it.
The F1 drop $\Delta_{100\to 5\%}$ measures how robust each method is
to attack-label scarcity.}
\small
\begin{tabular}{p{4.5cm}p{3.5cm}p{4cm}}
\toprule
\textbf{Configuration} & $F_1$ at 100\% labels & $\Delta F_1$ (100\%$\to$5\%)\\
\midrule
No pruning (Dense GCN) & $0.942\pm 0.005$ & $-0.210$ (std explodes)\\
\hline
$-$ learned, $+$ random (DropEdge) & $0.784\pm 0.006$ & $-0.047$ \\
\hline
$+$ learned, $-$ AE (NeuralSparse) & $0.661\pm 0.052$ & $-0.050$ \\
\hline
$+$ AE on \emph{full data} & $0.695\pm 0.067$ & $-0.003$ \\
\hline
\textbf{$+$ class-conditional} \textbf{(Ours)} & $\bm{0.743\pm 0.111}$
& $\bm{-0.001}$ \\
\bottomrule
\end{tabular}
\end{table}

Reading the table:
\begin{itemize}
  \item \textbf{Pruning is the trade-off}: dropping it gives the
        highest absolute $F_1$ (0.942) but at full graph cost and
        unstable behaviour at low labels.
  \item \textbf{Learned $>$ random}: ``learned but supervised''
        (NeuralSparse) is worse than ``learned but random'' (DropEdge)
        in this experiment, because the task gradient becomes a noisy
        signal under heavy graph noise.
  \item \textbf{Autoencoder gradient is the right signal}: adding the
        AE (even on full data) recovers absolute $F_1$ and reduces
        degradation slope.
  \item \textbf{Class-conditional masking completes the design}: the
        only change between row 4 and row 5 is the
        \texttt{(y == normal) \& train\_mask} filter on the recon
        loss, and that single line earns the flat curve.
\end{itemize}

% =========================================================================
\section{Connecting the math to the empirical results}
\label{sec:results-link}

\begin{table}[h]
\centering
\caption{Each empirical headline traces back to a specific piece of
math.}
\small
\renewcommand{\arraystretch}{1.4}
\begin{tabular}{p{6cm}p{8cm}}
\toprule
\textbf{Empirical result} & \textbf{Mathematical reason} \\
\midrule
Zero $F_1$ loss at 48\% edge reduction
(low noise, single seed)
& Pruner removes edges with low keep-score $s_{uv}$, which are exactly
the edges where $|\mathrm{MSE}_u - \mathrm{MSE}_v|$ is large
(cross-class). Removing them helps the GCN by reducing homophily
violation.\\
\hline
Selectivity ratio $1.42\times$ (vs $0.88\times$ for full AE)
& Under class-conditional training the AE learns
$p_X(\cdot\mid y=0)$, so
$\mathbb{E}[|\mathrm{MSE}_u-\mathrm{MSE}_v|\mid\text{cross-class}] \gg
\mathbb{E}[|\mathrm{MSE}_u-\mathrm{MSE}_v|\mid\text{same-class}]$.
Under full-AE training the AE learns the mixture and the gap collapses.\\
\hline
Flat $F_1$ curve from 100\% $\to$ 5\% attack labels
& The autoencoder's loss path
(Eq.~\eqref{eq:lrecon}) never references the attack label. Shrinking
$|L_{\text{train}}^{\text{attack}}|$ affects only $\mathcal{L}_{\text{CE}}$,
not the AE-derived OOD score.\\
\hline
$1.37\times$ inference speed-up at 50K-node graphs
& GCN forward cost is
$O(L\,|E|\,d + L\,|V|\,d^2)$.
Pruning 50\% of edges saves the first term. At small graphs the
second term dominates compute (so the speed-up is tiny); at large
graphs the first term dominates and the savings materialise.\\
\bottomrule
\end{tabular}
\end{table}

% =========================================================================
\section{The defence elevator-pitch}

\begin{keybox}[Memorise this 60-second answer]
\itshape
We start with the 3-layer GCN from Villegas-Ch et al.\ as the
classifier --- the architecture is unchanged. We add two components
on top: a class-conditional autoencoder that learns the normal-traffic
distribution and produces a per-node out-of-distribution score, and a
Gumbel-Softmax pruner with the Straight-Through estimator that uses
this OOD score together with cosine similarity and edge weight to
make differentiable, binary edge-keep/prune decisions. All three
components --- AE, pruner, and GCN --- are trained jointly with a
tri-partite loss whose three terms have three different supervision
modes: cross-entropy on labelled nodes (supervised),
MSE reconstruction on normal-class training nodes only
(self-supervised), and a target-sparsity quadratic penalty
(label-free). The methodological consequence is that the pruning
loss path uses no attack labels at all, so the pruning quality is
invariant to how many attack labels are available --- which is what
gives us the flat $F_1$ curve across the entire attack-label range,
and which closes the precise gap between supervised graph-pruning
methods (NeuralSparse, PTDNet) and unsupervised graph-anomaly
detection methods (DOMINANT, AnomalyDAE).
\end{keybox}

% =========================================================================
\section{One-page cheat-sheet}

\begin{tabular}{p{0.95\textwidth}}
\toprule
\textbf{The novelty in one sentence}\\
Train an autoencoder \emph{only on normal-class training nodes} so its
per-node MSE becomes an out-of-distribution score; feed this OOD score
to a differentiable Gumbel-Softmax pruner that runs before the
supervised GCN classifier.\\
\midrule
\textbf{The gap filled}\\
Existing pruners (NeuralSparse, PTDNet, DropEdge, GNNGuard) drive their
edge selection from the supervised classification loss $\Rightarrow$
they degrade when attack labels are scarce. Existing graph-anomaly
detectors (DOMINANT, AnomalyDAE) use AE reconstruction at the
\emph{node} level but never modify the graph. We sit at the
intersection: AE reconstruction $\Rightarrow$ \emph{edge-level}
pruning.\\
\midrule
\textbf{The three loss terms}\\
$\mathcal{L} \;=\; \mathcal{L}_{\text{CE}}
+ \lambda_1 \mathcal{L}_{\text{recon}}
+ \lambda_2 \mathcal{L}_{\text{sparse}}$
\quad with $\lambda_1{=}0.1$, $\lambda_2{=}1.0$.\\[2pt]
$\mathcal{L}_{\text{CE}}$: cross-entropy on $L_{\text{train}}$
(supervised).\\
$\mathcal{L}_{\text{recon}}$: MSE on $\{i\in L_{\text{train}}, y_i=0\}$
(self-supervised, no attack labels).\\
$\mathcal{L}_{\text{sparse}}$: target keep-rate constraint
(label-free).\\
\midrule
\textbf{Key empirical claims}\\
\begin{tabular}{p{0.45\textwidth}p{0.45\textwidth}}
$F_1$ = 0.995 at 48\% edge reduction (low noise) & vs Dense GCN 0.995\\
Selectivity ratio = $1.42\times$ & vs $0.88\times$ for full AE\\
$\Delta F_1$ from 100\% $\to$ 5\% labels = $-$0.001 & vs $-$0.210 for Dense GCN\\
Inference speed-up at 50K nodes = $1.37\times$ & deploy-mode\\
\end{tabular}\\
\midrule
\textbf{Anticipated panel questions \& one-line answers}\\
\textbf{``Why not NeuralSparse?''} It needs attack labels for its pruner
$\Rightarrow$ fails when labels scarce. Ours doesn't.\\
\textbf{``Is this just a regularization trick?''} No --- it's
\emph{supervision-decoupling}. Our pruner can train on data the
classifier cannot (unlabelled normals).\\
\textbf{``Why lower absolute $F_1$ at high noise?''} Trading accuracy for
50\% edge reduction, low variance under label scarcity, and inference
speed-up at scale. Among pruners we have highest mean \emph{and} lowest
variance at 5\% labels.\\
\textbf{``Why synthetic data?''} Need ground-truth edge-noise labels
for selectivity analysis. Real datasets don't have these. Future work.\\
\bottomrule
\end{tabular}

\vspace{1em}
\begin{notebox}[How this walkthrough relates to the formal thesis]
The thesis (\texttt{THESIS.tex}) is the formal academic write-up
intended for the panel. This walkthrough is the explainer intended
for \emph{you} --- to make sure you fully understand the math and can
explain every block. When defending, speak from the thesis; when
preparing, study this document.
\end{notebox}

\end{document}
